\begin{proposition}[Deformations of the constructed torus fibrations]\hypertarget{tag-01GW}{}\label{cdp:st:E7}
\StatementRecord{cdp:st:E7}{proved}{01GW}{\SeriesRef{cdp:st:E6}, \SeriesRef{cdp:eq:E-log-period}, \SeriesRef{cdp:st:E1}; \cite[Proposition~6.1]{Viaclovsky2026Combs}; \cite[Proposition~3.1]{Viaclovsky2026Two}; \cite[Proposition~2.5]{LiuRollenske2026}; \cite[Proposition~4.1, case~(i)]{Viaclovsky2026Combs}}

Every member of the following torus-fibred families,
including the period family $X_c$, its members with $|p|>1$,
Viaclovsky's $X_{3k,4l}$, his two-parameter family and his combs,
and the Liu--Rollenske family, satisfies the rank hypothesis of \ref{cdp:st:E6}.
Consequently every sufficiently small deformation retains a
holomorphic torus fibration over $\mathbb P^1$. Every such member
with $b_2=0$, and every sufficiently small deformation of it, has
algebraic dimension one. In particular the fibration-retention
question in \cite[Section~10(1)]{Viaclovsky2026Combs} has an
affirmative answer.
\end{proposition}